\documentclass[11pt]{letter}
\usepackage[margin=0.85in]{geometry}
\usepackage[T1]{fontenc}
\usepackage{lmodern}
\usepackage[hidelinks]{hyperref}
\usepackage{microtype}

\pagestyle{empty}
\setlength{\parindent}{0pt}
\setlength{\parskip}{0.8em}

\signature{Adam Bawatneh}
\address{Adam Bawatneh\\Orlando, Florida\\(718) 781-7009\\\href{mailto:adambawatneh01@gmail.com}{adambawatneh01@gmail.com}\\\href{https://adam-12-0.github.io/}{adam-12-0.github.io}}

\begin{document}
\begin{letter}{Hiring Team\\Databricks\\San Francisco, California}
\date{October 16, 2026}
\opening{Dear Databricks Hiring Team,}

I am a Ph.D. student in Computer Science at the University of Central Florida, expecting to graduate in Spring 2029, and I am applying for the Ph.D. GenAI Research Scientist Intern position. My research focuses on social reasoning and Theory of Mind, with a broader interest in how language models represent information and generalize across contexts. I am drawn to this role because it combines careful evaluation of AI systems with research on adapting them to new domains.

I led OmniToM, a paper accepted to the NeurIPS 2026 Evaluations \& Datasets Track on evaluating Theory of Mind in large language models. I developed its 895-story benchmark and coordinated quality review with 11 annotators across more than 1,100 hours. I also evaluated commercial and open models through cloud APIs. This work gave me experience designing evaluations, organizing human-labeled data, and analyzing how model behavior changes across challenging scenarios.

My work on CADS, a project on detecting specialized language, has further developed my NLP research skills. I built and evaluated a text-processing pipeline and examined how its performance transferred across domains and languages. That experience has made me interested in methods that help AI systems work reliably with data from different settings, which closely connects to Databricks' research on domain adaptation, language-model evaluation, and AI systems for enterprise data.

I would bring a careful experimental approach, experience building research datasets and evaluation pipelines, and the ability to communicate findings to technical collaborators. I would welcome the opportunity to contribute these strengths to the Databricks research team and learn from its work on practical generative AI systems.

\closing{Sincerely,}
\end{letter}
\end{document}
